\documentclass[12pt]{article}
\usepackage[utf8]{inputenc}
\usepackage[T1]{fontenc}
\usepackage{hyperref}
\title{Punan Merah Language Resources for K-12 Curriculum Development}
\author{Compiled by Claude Sonnet 5.5 (claude-sonnet-5-5) for Andr\'as Kornai}
\date{}
\begin{document}
\maketitle

\section{Introduction}
Punan Merah (ISO 639-3 \texttt{puf}, Glottocode puna1268) is an Austronesian language of East Kalimantan, Indonesia, spoken by the Merah (Red) Punan, who live in East and North Kalimantan. The only speaker figure found is about 140 speakers, from 1981, so it is old. It is listed as endangered in a survey of Kalimantan languages. Searches found no dedicated grammar, dictionary, primer, scripture or school materials.

\section{Primary Educational Resources}
None found.

\section{Linguistic and Reference Materials}
\begin{itemize}
\item Punan Merah language, Wikipedia (English): overview and speaker estimate. \url{https://en.wikipedia.org/wiki/Punan_Merah_language}
\item Soriente and Inagaki, 2012 overview slides on Kalimantan languages; notes limited documentation and an unresolved Punan classification. \url{https://lingdy.aa-ken.jp/wp-content/uploads/2012/01/120217_soriente_inagaki_h2.pdf}
\end{itemize}

\section{Religious and Community Resources}
\begin{itemize}
\item Merah Punan in Indonesia, Joshua Project: people-group profile. \url{https://joshuaproject.net/people_groups/14479}
\end{itemize}

\section{Digital and Audio Resources}
None found.

\section{Curriculum Development Context}
There is almost no published documentation, so a curriculum would have to begin with community-led documentation and an agreed orthography. Indonesian remains the language of schooling. The sources above are general reference pages, not teaching materials.

\end{document}